% !Mode:: "TeX:UTF-8"
% 北京工业大学数据挖掘课程报告宏包配置
\usepackage{iftex}
\ifPDFTeX
  \PackageError{BJUTCourseReport}
    {This template must be compiled with XeLaTeX or LuaLaTeX}
    {In Overleaf, open Menu -> Compiler and choose XeLaTeX.}
\fi
% ============================================================
% 1. 页面布局与尺寸 (核心：适配双面打印与36行网格)
% ============================================================
\usepackage[
    a4paper,
    top=25mm,
    bottom=25mm,
    inner=27mm,    % 装订内侧边距
    outer=27mm,    % 外侧边距
    headheight=5mm,
    headsep=7.5mm, % 页眉位置
    footskip=12mm, % 页脚位置
    lines=36,      % 显式指定行数 (配合 format.tex 中的字号设置)
    heightrounded, % 自动微调高度以容纳整数行
]{geometry}

% ============================================================
% 2. 数学公式支持
% ============================================================
\usepackage{amsmath}   % 基础数学宏包
\usepackage{amssymb}   % 数学符号
\usepackage{amsfonts}  % 数学字体
\usepackage{bm}        % 数学粗体
\usepackage[amsmath,thmmarks,hyperref]{ntheorem} % 定理环境

% ============================================================
% 3. 图表支持
% ============================================================
\usepackage{graphicx}  % 插图核心宏包
\graphicspath{{figures/}} % 图片默认路径

\usepackage{booktabs}  % 三线表 (\toprule, \midrule, \bottomrule)
\usepackage{multirow}  % 表格多行合并
\usepackage{longtable} % 跨页长表格
\usepackage{array}     % 增强表格功能
\usepackage{caption}   % 图表标题定制 (替代老旧的 ccaption)
\usepackage[below]{placeins} % 防止浮动体跨过 \FloatBarrier
\usepackage{bicaption} % 支持双语图注
% ============================================================
% 4. 格式控制与辅助工具
% ============================================================
\usepackage{titlesec}  % 章节标题格式控制
\usepackage{titletoc}  % 目录格式控制
\usepackage{fancyhdr}  % 页眉页脚控制
\usepackage{enumitem}  % 列表环境控制
\usepackage{setspace}  % 行距控制
\usepackage{xcolor}    % 颜色支持
\usepackage{fontspec}  % 字体高级支持 (配合 XeLaTeX)
\usepackage{calc}      % 长度计算

% ============================================================
% 5. 参考文献
% ============================================================
\usepackage[sort&compress,numbers,super]{natbib} % 引用优化 (排序+压缩: [1-3])

% ============================================================
% 6. 代码高亮
% ============================================================
\usepackage{listings}

% 定义颜色
\definecolor{codegreen}{rgb}{0,0.6,0}
\definecolor{codegray}{rgb}{0.5,0.5,0.5}
\definecolor{codepurple}{rgb}{0.58,0,0.82}
\definecolor{backcolour}{rgb}{0.95,0.95,0.92} % 护眼淡灰色背景

% 设置代码样式
\lstdefinestyle{mystyle}{
    backgroundcolor=\color{backcolour},
    commentstyle=\color{codegreen}\itshape, % 注释斜体+绿色
    keywordstyle=\color{blue}\bfseries,     % 关键字蓝色+加粗
    numberstyle=\tiny\color{codegray},      % 行号样式
    stringstyle=\color{codepurple},         % 字符串样式
    basicstyle=\zihao{5}\ttfamily,          % 字号五号 + 等宽字体
    breakatwhitespace=false,
    breaklines=true,                 % 自动换行
    captionpos=b,                    % 标题在底部
    keepspaces=true,
    numbers=left,                    % 行号在左侧
    numbersep=5pt,
    showspaces=false,
    showstringspaces=false,
    showtabs=false,
    tabsize=4,
    frame=shadowbox,                 % 阴影边框 (也可以选 single)
    rulesepcolor=\color{red!20!green!20!blue!20}, % 阴影颜色
    xleftmargin=2em,                 % 左侧缩进，防止行号溢出
    xrightmargin=1em,                % 右侧缩进
    framexleftmargin=1.5em           % 边框左侧扩充
}

\lstset{style=mystyle}

% ============================================================
% 6.5. 伪代码算法环境
% ============================================================
\usepackage[ruled,vlined,linesnumbered]{algorithm2e}

% ============================================================
% 8. 字体设置：适配 Overleaf / TeX Live
% ============================================================
\usepackage{fontspec}

\ifPDFTeX\else
  \setmainfont{TeX Gyre Termes}
  \setsansfont{TeX Gyre Heros}
  \setmonofont{TeX Gyre Cursor}

  \newcommand{\BJUTSetCJKFamilies}[3]{%
    \setCJKmainfont{#1}%
    \setCJKsansfont{#2}%
    \setCJKmonofont{#2}%
    \setCJKfamilyfont{zhsong}{#1}%
    \setCJKfamilyfont{zhhei}{#2}%
    \setCJKfamilyfont{zhkai}{#3}%
    \providecommand{\songti}{}%
    \providecommand{\heiti}{}%
    \providecommand{\kaishu}{}%
    \providecommand{\fangsong}{}%
    \renewcommand{\songti}{\CJKfamily{zhsong}}%
    \renewcommand{\heiti}{\CJKfamily{zhhei}}%
    \renewcommand{\kaishu}{\CJKfamily{zhkai}}%
    \renewcommand{\fangsong}{\CJKfamily{zhsong}}%
  }

  \IfFontExistsTF{Noto Serif CJK SC}{%
    \BJUTSetCJKFamilies{Noto Serif CJK SC}{Noto Sans CJK SC}{Noto Serif CJK SC}%
    \IfFontExistsTF{Noto Sans Mono CJK SC}{\setCJKmonofont{Noto Sans Mono CJK SC}}{}%
  }{%
    \IfFontExistsTF{FandolSong-Regular}{%
      \BJUTSetCJKFamilies{FandolSong-Regular}{FandolHei-Regular}{FandolKai-Regular}%
      \setCJKmonofont{FandolFang-Regular}%
    }{%
      \BJUTSetCJKFamilies{AR PL UMing CN}{AR PL UMing CN}{AR PL UMing CN}%
    }%
  }%
\fi

% ============================================================
% 7. 超链接设置
% ============================================================
\usepackage[
    unicode=true,
    bookmarksnumbered=true,
    bookmarksopen=true,
    colorlinks=false, % 打印版通常要求黑色链接
    pdfborder={0 0 0}, % 去掉链接红框
    pdfkeywords={},
    pdfstartview=FitH
]{hyperref}

% ============================================================
% 生僻字显示方案：字体已在上方设置，通常无需额外处理。
\endinput
